% example_book.tex — demonstrates colegio.cls in book mode
%
% A miniature two-chapter book demonstrating: title page with full
% protocol metadata, \chapter-based sectioning, tone-aware chapter rule,
% all three annotation modes within chapter prose, multi-chapter \backnote
% collection at the end, and the colophon environment with inscription
% details + unattached annotations surfaced.
%
% Compile: xelatex example_book.tex && xelatex example_book.tex
\documentclass[book,tone=ordinary]{../colegio}

\title{Two Sample Chapters}
\persona{El Gólem}
\date{2026-05-23}
\blockheight{6{,}218{,}xxx}
\inscribingaddress{example address (illustrative only)}
\rootxid{0000000000000000000000000000000000000000000000000000000000000000}
\jointxid{0000000000000000000000000000000000000000000000000000000000000000}
\protocolheader{c1 dd 00 01 \,09 00\, |Two Sample Chapters|author=El G\'olem|date=2026-05-23|lang=en|}

\begin{document}

\maketitle

\frontmatter
\tableofcontents

\mainmatter

\chapter{On the Cord}

\newthought{A book is}, in the protocol, a 0x09 inscription whose body lists
references to other quipus. Each reference carries a tag, a name, and
a 32-byte ref\_txid. The book is not the prose; the book is the
manifest.\margin{The same essay rendered through two different books
can present differently --- each book's binding entries provide a
contextual overlay (title rewrites, citations, annotations) without
altering the inscribed prose.}

The cord runs from the root transaction through the strands; the knots
on each strand are the OP\_RETURN-bearing transactions; the join
collects every strand's terminus into one output that returns the
unused capital to the inscriber's address. This is the
\emph{diamond}.\backnote{the diamond}{Each diamond is one inscription:
a splitter transaction with N outputs (one per strand), N parallel
strand chains, and a single mega-join transaction collecting all the
strand termini. The geometry is dictated by Dogecoin's mempool
ancestor limit of 25, which caps how long a single strand can be
before requiring a block-confirmation wait.}

A book that points at another book is a library\inlinenote{a book of
books, recursive by virtue of the type system}. The recursion is
bounded by the reader's chosen depth limit; the project's default is
eight levels, matching the binding alias-chain depth.

\chapter{On the Knot}

\newthought{Every knot} is one transaction. Every transaction carries one
OP\_RETURN. Every OP\_RETURN carries up to eighty bytes of payload.
\margin{eighty bytes}{The Bitcoin / Dogecoin script primitive
\texttt{OP\_RETURN} provably-burns the output it appears in, signaling
to validators that the eighty bytes following are data not value. The
fee for emitting them is the same as any other transaction's fee:
small.} The total payload of an inscription is therefore the sum of
its strand-knot OP\_RETURNs, concatenated in strand order, parsed by
the canonical type module corresponding to the inscription's type
byte.

Reading a quipu is following the cord. The reader fetches the root
transaction, identifies its strand-seed outputs, walks each strand
forward through its OP\_RETURN-bearing knots, and concatenates the
payloads. There is no global index; there is no central catalog.
The chain is the catalog, and the catalog is the chain.

Citations of other quipus are written as
\quipucite[Dos ensayos]{26671514416913719c560a4ac1246e333e410d2503e298a8a6852031c3285888}
in the inscribed prose. The reader resolves the citation by fetching
the cited txid, parsing its header, and rendering or linking
accordingly.

\backmatter

\printbacknotes

\begin{colophon}
  \colophonentry{author}{El Gólem}
  \colophonentry{date}{2026-05-23}
  \colophonentry{lang}{English}
  \colophonentry{type byte}{0x09 book}
  \colophonentry{tone byte}{0x00 ordinary}
  \colophonentry{address}{(example address)}
  \colophonentry{root txid}{0000000000000000000000000000000000000000000000000000000000000000}
  \colophonentry{rendering}{xelatex via \texttt{colegio.cls} v0.1 (book mode)}
  \colophonentry{note}{This is a class-demonstration document. The txids above are placeholders and not inscribed on chain.}
  % Demonstrate one unattached annotation (anchor not present in the body):
  \unattached{a phrase that doesn't appear in the prose}{This is what
  the unattached-annotation mechanism does: the inscriber's intent is
  preserved even when the anchor mis-matches. The note surfaces here
  in the colophon rather than being silently dropped.}
\end{colophon}

\end{document}
